


\section{Limitations and Discussion}
The main limitation of our approach is that the base models need to have sufficient few-shot modeling or instruction-following abilities, in order to learn to provide feedback and to refine in an in-context fashion, without having to train supervised models and rely on supervised data.

Further, the experiments in this work were performed with language models that are not open-sourced, namely \gptt, \chatgpt, \gptf, and \codex. 
Existing literature~\citep{Ouyang2022TrainingLM} does not fully describe the details of these models, such as the pretraining corpus, model sizes, and model biases.
Further, these models are not free to use, and using them for research requires some funding.
Nonetheless, we release our code and model outputs to ensure the reproducibility of our work. 


Another limitation of our work is that we exclusively experiment with datasets in English. In other languages, the current models may not provide the same benefits.

Finally, there is a possibility for bad actors to use prompting techniques to steer a model to generate more toxic or harmful text. Our approach does not explicitly guard against this.


\section{Conclusion}





We present \ours: a novel approach that allows large language models to  iteratively provide self-feedback and refine their own outputs.
\ours %
operates within a single \llm, requiring neither additional training data nor reinforcement learning. 
We demonstrate the simplicity and ease of use of \ours 
across a wide variety of tasks. %
By showcasing the potential of \ours in diverse tasks, our research contributes to the ongoing exploration and development of large language models, with the aim of reducing the cost of human creative processes in real-world settings.
We hope that our iterative approach %
will help drive further research in this area.
To this end, we make all our code, data and prompts anonymously available at 
\url{https://selfrefine.info/}. 

